%! Logarithmic Expressions — Unit 5, Lesson 5.1 — Student Packet Cover Sheet
\documentclass[10pt]{article}
\usepackage{precalculus-article}
\usepackage{precalculus-boxes}
\usepackage{ltablex}
\keepXColumns

\begin{document}

% Full-bleed navy banner
\begin{tikzpicture}[remember picture, overlay]
  \fill[navy] (current page.north west)
    rectangle ([yshift=-0.9in]current page.north east);
\end{tikzpicture}
\vspace{-0.8in}

\begin{center}
  {\color{white}\textbf{\LARGE Precalculus}}\\[4pt]
  {\color{skymid}\large\textbf{Unit 5: Logarithmic Functions}}\\[6pt]
  {\color{white}\textbf{\large Lesson 5.1 \quad Logarithmic Expressions}}
\end{center}

\vspace{0.22in}
\namedateperiod
\vspace{0.18in}

\begin{learningtargetbox}
\small
\begin{enumerate}[leftmargin=1.5em, itemsep=5pt]
  \item I can convert between exponential form ($b^a = c$) and logarithmic form
        ($\log_b c = a$) and evaluate exact logarithms.
  \item I can evaluate common logarithms ($\log_{10}$) and natural logarithms ($\ln$),
        and use technology to estimate logarithm values.
  \item I can interpret a logarithmic scale, recognizing that each unit represents a
        multiplicative change by the base.
\end{enumerate}
\end{learningtargetbox}

\vspace{0.14in}

\begin{tocbox}
\small
\renewcommand{\arraystretch}{1.5}
\begin{tabularx}{\linewidth}{@{} c l X r @{}}
\rowcolor{sky}
\textbf{\#} & \textbf{Component} & \textbf{Description} & \textbf{Score} \\
\midrule
1 & Warm-Up        & Spiral review: inverse functions and exponent identification & \blank{1.2cm} \\
2 & Guided Notes   & Definition, conversion, common/natural log, logarithmic scale & \blank{1.2cm} \\
3 & Group Activity & Conversion drill and applications (tiered R/A/E) & \blank{1.2cm} \\
4 & Exit Ticket    & Convert, evaluate, and bound a logarithm & \blank{1.2cm} \\
5 & Homework       & Mixed: conversion, exact evaluation, technology, real-world, AP MC & \blank{1.2cm} \\
\midrule
  & \multicolumn{2}{r}{\textbf{Total}} & \blank{1.2cm} \\
\end{tabularx}
\end{tocbox}

\end{document}
